\begin{abstract}
La Direction Générale de la Sûreté Nationale (DGSN) délivre chaque année plusieurs millions de cartes nationales d'identité électroniques (CNIE) et opère, depuis 2022, la plateforme nationale d'identité numérique. Le portail \texttt{cnie.ma} dématérialise la pré-demande et la prise de rendez-vous, mais l'essentiel du contrôle des dossiers reste manuel et les canaux numériques de la DGSN demeurent cloisonnés. Ce rapport applique la démarche d'urbanisation des systèmes d'information au périmètre «~demande et renouvellement de la CNIE~». Il capte d'abord les orientations stratégiques de l'organisation (Digital Morocco 2030, loi 55-19, lutte contre la fraude, inclusion), puis décline le système cible en quatre cartographies cohérentes~: métier (processus principal modélisé en BPMN sous IBM Blueworks Live, dans ses versions actuelle et cible), fonctionnelle (plan d'occupation des sols distinguant blocs existants, à faire évoluer et nouveaux), applicative (quinze composants cibles et leurs flux) et technique (architecture en zones, cloud souverain, PKI, plateforme d'inférence). Cinq cas d'usage d'intelligence artificielle sont intégrés au cœur du processus~: assistant conversationnel, contrôle documentaire automatique, vérification biométrique anti-fraude, prévision d'affluence et détection d'anomalies. Le sous-processus de pré-demande, de contrôle des pièces et de prise de rendez-vous a été prototypé sous Bonitasoft pour valider l'orchestration et les interfaces. Une trajectoire en trois paliers (2027, 2028, 2030) et un dispositif de gouvernance complètent la proposition.
\end{abstract}

\begin{IEEEkeywords}
urbanisation des SI, architecture d'entreprise, ArchiMate, BPMN, identité numérique, CNIE, DGSN, intelligence artificielle, e-gouvernement, Bonita, Blueworks Live
\end{IEEEkeywords}
